\section{Introduction}
\label{sec:introduction}

Large language models (LLMs) have enabled a new class of applications that solve tasks not by a single inference call, but by executing \emph{agentic workflows}: structured compositions of multiple LLM invocations, tool calls, and control logic such as branching, voting, and iterative refinement. Frameworks such as MetaGPT~\citep{hong2023metagpt} and AutoGen~\citep{wu2023autogen} show that decomposing a task into role-specialized agents can improve robustness, while systems such as HuggingGPT~\citep{shen2023hugginggpt} demonstrate that LLMs can act as controllers to orchestrate heterogeneous expert models and tools. More recently, graph-centric workflow systems (e.g., GPTSwarm~\citep{zhuge2024gptswarm} and LangGraph~\citep{langgraph2024docs}) make workflow structure explicit and programmable, enabling automated topology search~\citep{zhang2024aflow,zhang2025maas} and empirically grounded studies of coordination regimes~\citep{kim2025scalingagents}.

Despite rapid progress in agent construction, the end goal for deployment is increasingly a systems problem: agentic workflows amplify prefill redundancy, induce bursty parallelism, introduce synchronization barriers, and interleave GPU-bound LLM execution with unbounded tool latency. These behaviors stress the serving stack in ways that are invisible under the conventional ``one request = one generation'' abstraction. Recent serving systems have begun to expose application structure to the backend---e.g., ParrotServe's semantic variables~\citep{lin2024parrot} and KVFlow's workflow-aware KV management~\citep{pan2025kvflow}---but the broader design space remains fragmented: workflow papers propose diverse abstractions (graphs, code, conversations), while system papers optimize isolated motifs (prefix caching, scheduling, disaggregation) without a unifying formal model of agent execution.

This survey argues that \textbf{formalizing agentic workflows as a well-defined graph abstraction} is the missing bridge between frontend agent design and backend system optimization. Our central motivation is to make workflow structure structurally visible so that it becomes analyzable, comparable across frameworks, and optimizable by the runtime.

\paragraph{What we do.}
We (1)~survey representative agentic workflow frameworks and extract their underlying execution structure; (2)~formalize these workflows under a unified typed graph abstraction $G=(\mathcal{N},\mathcal{E},\mathrm{Op})$ that treats computation, tools, data dependencies, and control operators as first-class entities; and (3)~use this abstraction to systematically organize system optimization opportunities. Concretely, we show how recurring topological motifs (e.g., fan-out, fan-in, sequential chains, self-cycles, star routing, and compute--I/O alternation) map to execution-level optimizations (scheduling, KV reuse, prefetch/eviction, placement), and how graph-level transformations (early exits, node downscaling, router--expert decomposition, subgraph consolidation, and topology sparsification) reshape workflows to reduce total inference cost.

\paragraph{Why this matters.}
A unified abstraction turns ad-hoc workflow engineering into a structured optimization problem: given a query distribution and system constraints (latency SLOs, GPU memory, tool overhead), the runtime can reason about alternative workflow graphs and select the best cost--quality tradeoff. This perspective also enables the next step---graph auto-optimization---where the workflow is treated as an optimizable object (e.g., probabilistic edges~\citep{zhuge2024gptswarm} or topology search~\citep{zhang2024aflow}) but with system efficiency as a first-class objective.

\section{Related Work}
\label{sec:related_work}

\paragraph{LLM Serving Optimizations.}
A broad line of work improves online LLM serving via better scheduling and memory management. vLLM introduces PagedAttention to reduce KV fragmentation and enable flexible KV allocation and sharing~\citep{kwon2023efficient}. SGLang exploits prefix reuse by organizing KV segments in a radix tree (RadixAttention) and performing efficient longest-prefix matching~\citep{zheng2024sglang}. Sarathi-Serve mitigates throughput--latency interference by chunking large prefills and interleaving them with ongoing decode, enabling stall-free schedules under bursty prefill workloads~\citep{agrawal2024taming}. Beyond caching and scheduling, disaggregated architectures split prefill and decode onto different resources to reduce phase interference and improve utilization, as exemplified by DistServe~\citep{zhong2024distserve} and Splitwise~\citep{patel2024splitwise}. ServerlessLLM targets multi-model deployments by reducing cold-start and checkpoint loading overheads via locality-aware model storage~\citep{fu2024serverlessllm}. More recently, programmable serving systems such as Pie expose fine-grained control over KV, generation, and I/O to application logic, enabling workflow-specific optimizations beyond generic prefix caching~\citep{gim2025pie}.

\paragraph{Serving and Optimization for Agentic Workflows.}
Agentic workloads introduce structured multi-step execution, tool interleaving, and synchronization barriers that conventional serving stacks do not explicitly model. ParrotServe exposes prompt-level dependencies via semantic variables and uses DAG-aware scheduling and context forking to reduce redundant computation across chained LLM calls~\citep{lin2024parrot}. KVFlow focuses specifically on agentic workflows and proposes a workflow-aware cache eviction and fully overlapped KV prefetching based on an Agent Step Graph abstraction~\citep{pan2025kvflow}. Complementary directions study batch-level processing and optimization for agentic workflows (e.g., integrating adaptive batching, KV sharing/migration, and plan-level optimizations)~\citep{halo2025}. These systems collectively support our thesis that workflow structure must be surfaced to the backend to unlock principled systems optimizations.

\paragraph{Agentic Workflow Frameworks and Graph-Based Orchestration.}
A large body of work proposes reasoning-and-acting workflows (e.g., ReAct, Reflexion, Tree-of-Thought) and multi-agent coordination patterns (e.g., debate and role-based collaboration), many of which can be expressed as graphs of LLM calls and tools~\citep{zhuge2024gptswarm}. At the framework level, MetaGPT encodes standardized operating procedures as multi-role pipelines to improve coordination in software engineering tasks~\citep{hong2023metagpt}, while AutoGen provides ``conversable agents'' and conversation programming to compose multi-agent applications with tools and optional human inputs~\citep{wu2023autogen}. HuggingGPT uses an LLM controller to plan, select, and orchestrate heterogeneous expert models and tools to solve multimodal tasks~\citep{shen2023hugginggpt}. LangGraph provides a low-level runtime for long-running, stateful workflows with persistence, interrupts, and subgraphs, exposing orchestration structure directly to developers~\citep{langgraph2024docs}. GPTSwarm formalizes agents as DAGs and swarms as composite ``graphs of graphs,'' enabling both node-level prompt refinement and edge-level topology optimization via reinforcement learning~\citep{zhuge2024gptswarm}.

\paragraph{Automatic Workflow Design, Topology Search, and Adaptation.}
Several recent works treat workflow structure as a search or adaptation object. AFlow formulates workflow optimization as MCTS over workflows with LLM-invoking nodes, code-represented control edges, and reusable operators (e.g., ensemble, review--revise)~\citep{zhang2024aflow}. MaAS introduces an agentic supernet that samples query-dependent subnetworks to allocate computation adaptively~\citep{zhang2025maas}. DAAO adapts workflow depth, operator choice, and LLM assignment based on difficulty estimation and cost-aware routing~\citep{su2025daao}. These approaches are closely aligned with our graph-level optimization view: they mutate or select subgraphs, but are typically evaluated primarily on accuracy or monetary cost; our emphasis is to connect such transformations to backend objectives (latency, throughput, KV residency, and I/O overlap) under a unified formal model.

\paragraph{Scaling Laws, Topology Tradeoffs, and Consolidation.}
Understanding when coordination helps remains an open question. ``Towards a Science of Scaling Agent Systems'' performs controlled evaluations over canonical coordination topologies (single, independent, centralized, decentralized, hybrid) and derives predictive principles based on measurable coordination metrics~\citep{kim2025scalingagents}. In parallel, recent work questions the necessity of homogeneous multi-agent workflows: OneFlow shows that many homogeneous multi-agent workflows can be simulated by a single agent via multi-turn execution, benefiting from KV reuse within one model context~\citep{xu2026oneflow}. AgentArk further distills multi-agent debate dynamics into a single model's weights, shifting coordination cost from inference-time orchestration to training-time internalization~\citep{luo2026agentark}. These results delineate boundary conditions for graph-level consolidation and motivate hybrid runtimes that choose between distributed and consolidated subgraphs based on predicted system cost and quality.